% ===================================================================
%  CHART No. 6 — The Home inside. — Furniture. — Food. — Lighting.
%  Traduction 2026 d'apres texto/io, controlee sur texto/fr.
%  Consigne : voir texto/en/10-table-01.tex.
%
%  Deux notes, toutes deux marquees « (*) », et deux appels : celui de
%  « babo » au bloc c1-12-1, celui de « lambrequino » au bloc c2-01-1.
%  L'ordre les departage, comme dans les deux autres colonnes.
%
%  LA FEMME DE CHAMBRE N'A PAS DE GROS PLAN, et c'est normal. L'ido du
%  bloc c1-06-1 porte « la chambristino (6) », le francais « la femme
%  de chambre (16) » : l'un des deux ateliers s'est trompe, et c'est
%  l'ido, puisque son propre numero 6 est la sponge du bloc c1-02-1.
%  L'anglais suit donc le francais et ecrit (16). Mais la planche ne
%  porte pas de 16 releve -- les numeros attendus se lisent sur le
%  texte IDO, qui n'en a jamais demande -- et le renvoi reste du texte
%  ordinaire, faute de savoir ou pointer. C'est la regle de la maison :
%  on ne promet pas un gros plan qu'on ne saurait pas montrer.
% ===================================================================

%%K t06-apar-1 apar
\VUsaut{6.0mm}
\VUtitre{15.0pt}{60}{CHART No. 6}
\VUsaut{1.4mm}
\VUtitre{11.4pt}{0}{The Home inside, Furniture, Food,}
\VUsaut{0.4mm}
\VUtitre{11.4pt}{0}{Lighting.}
\VUsaut{2.2mm}

%%K t06-apar-2 apar
The house is finished. John's uncle has moved into his new home, whose
layout we already know. Let us now see how he has furnished it. First we
shall visit:

%%K t06-tit-1 sub
\VUpk{10.2pt}{40}{The Bedroom.}
\VUsaut{3.0mm}

%%K t06-c1-01-1 p
1. --- We have come at an awkward moment, for the maid is busy cleaning
the \VUgras{room}.

%%K t06-c1-02-1 p
2. --- On the \VUgras{washstand} \textsuperscript{(1)} lie, here and there, the
various things one washes, combs and generally gets ready with. They are
the \VUgras{basin} \textsuperscript{(2)} and the \VUgras{water jug} \textsuperscript{(3)},
the \VUgras{hairbrush} \textsuperscript{(4)}, the \VUgras{comb} \textsuperscript{(5)}, the
\VUgras{soap} \textsuperscript{(6)} and the \VUgras{sponge} \textsuperscript{(7)}, and lastly a
little \VUgras{bottle} \textsuperscript{(8)} of mouthwash.

%%K t06-c1-03-1 p
3. --- On the floor, in front of the washstand, here is the
\VUgras{pail} \textsuperscript{(9)} for the dirty water.

%%K t06-c1-04-1 p
4. --- In the \VUgras{outer room} \textsuperscript{(10)} we can see a few
\VUgras{coat hooks} \textsuperscript{(13)} and two \VUgras{umbrellas} \textsuperscript{(11)}
in the \VUgras{umbrella stand} \textsuperscript{(12)}.

%%K t06-c1-05-1 p
5. --- On the other side of the door stands the \VUgras{mirrored wardrobe}
\textsuperscript{(14)}, which holds the \VUgras{linen} \textsuperscript{(15)}.

%%K t06-c1-06-1 p
6. --- Let us take advantage of the \VUgras{housemaid} \textsuperscript{(16)} making
the \VUgras{bed} \textsuperscript{(17)} to look at its various parts. The
\VUgras{bedstead} \textsuperscript{(20)} holds a good \VUgras{spring mattress}
\textsuperscript{(21)}, on which Justine will presently lay a woollen
\VUgras{mattress} \textsuperscript{(22)}. Then she will take from the chairs the two
\VUgras{sheets} \textsuperscript{(23)} and the \VUgras{blanket} \textsuperscript{(24)}. Then
she will place at the head of the bed the \VUgras{bolster} \textsuperscript{(25)} and
the \VUgras{pillow} \textsuperscript{(26)}, which are lying with the \VUgras{eiderdown}
\textsuperscript{(27)} on the \VUgras{bedside rug} \textsuperscript{(32)}. She uses a
feather duster on the \VUgras{curtains} \textsuperscript{(19)} and the
\VUgras{canopy} \textsuperscript{(18)}, and part of the work is done.

%%K t06-c1-07-1 p
7. --- Then she will have to tidy the \VUgras{bedside table} \textsuperscript{(28)} a
little, wind up the \VUgras{alarm clock} \textsuperscript{(29)}, get the
\VUgras{night light} \textsuperscript{(30)} ready, and take away the \VUgras{candle}
\textsuperscript{(31)} to clean the candlestick.

%%K t06-tit-2 sub
\VUsaut{5.0mm}
\VUpk{10.2pt}{40}{The workbasket and the fireside things.}
\VUsaut{3.0mm}

%%K t06-c1-08-1 p
8. --- The \VUgras{workbasket} \textsuperscript{(34)} beside the \VUgras{stool}
\textsuperscript{(33)} badly needs tidying too. She will have to fold up the
\VUgras{embroidery} \textsuperscript{(35)}, put away in the \VUgras{worktable}
\textsuperscript{(38)} the \VUgras{thimble}, the \VUgras{thread} \textsuperscript{(36)}, the
\VUgras{needles} \textsuperscript{(37)} and the \VUgras{scissors} \textsuperscript{(39)}, and
lock the lid of the table with the \VUgras{key} \textsuperscript{(40)}.

%%K t06-c1-09-1 p
9. --- On the \VUgras{pedestal table} \textsuperscript{(41)} in the middle, Justine
will change the water in the \VUgras{carafe} \textsuperscript{(42)}. She will put
\VUgras{sugar} in the \VUgras{sugar bowl} \textsuperscript{(43)}, clean the
\VUgras{sugar tongs} \textsuperscript{(44)} and take away the \VUgras{clothes brush}
\textsuperscript{(46)}.

%%K t06-c1-10-1 p
10. --- The right-hand side of the room will give her less work, for the
\VUgras{day bed} \textsuperscript{(47)} and its \VUgras{cushion} \textsuperscript{(48)} are
in their proper place, and since the weather is not cold, nothing has
been moved among the things belonging to the \VUgras{fireplace}
\textsuperscript{(49)}. \VUgras{Shovel} \textsuperscript{(50)}, \VUgras{tongs}
\textsuperscript{(51)}, \VUgras{firedogs} \textsuperscript{(55)} and \VUgras{ash guard}
\textsuperscript{(56)} are clean; the \VUgras{fire screen} \textsuperscript{(54)} has not
been moved and the \VUgras{log box} \textsuperscript{(52)} is well stocked with
\VUgras{firewood} \textsuperscript{(53)}.

%%K t06-noto-1 noto
\VUnotes{143.2mm}{%
(*) Babo = a small child; English \textit{baby}, French \textit{bébé},
Italian \textit{bambino}.%
}

%%K t06-noto-2 noto
\VUnotes{143.2mm}{%
(*) Lambrequino = French \textit{lambrequin}, Spanish
\textit{lambrequines}, Portuguese \textit{lambrequins}, and even German
and English \textit{lambrequins} --- so German, English, French,
Portuguese and Spanish alike.%
}

%%K t06-c1-11-1 p
11. --- She will still have to dust the \VUgras{mantel clock} \textsuperscript{(57)},
the \VUgras{candelabra} \textsuperscript{(58)} and the \VUgras{trinket trays}
\textsuperscript{(59)}, and then wipe the \VUgras{mirror} \textsuperscript{(60)}.

%%K t06-c1-12-1 p
12. --- As for the baby's \VUgras{cradle} \textsuperscript{(61)} (the babo (*)), it
is ready already.

%%K t06-tit-3 sub
\VUsaut{5.0mm}
\VUpk{10.2pt}{40}{The Drawing Room.}
\VUsaut{3.0mm}

%%K t06-c2-01-1 p
1. --- Here now is the drawing room. It is a very fine \VUgras{room}. The
fireplace, in the right-hand corner, is decorated with a delicate
\VUgras{statuette} \textsuperscript{(1)} and two handsome \VUgras{vases}
\textsuperscript{(3)}, and draped with a velvet \VUgras{lambrequin} \textsuperscript{(2)}
(*).

%%K t06-c2-02-1 p
2. --- To keep the sparks from the hearth from setting the carpet alight,
a copper \VUgras{fire guard} \textsuperscript{(4)} has been set in front of it.

%%K t06-c2-03-1 p
3. --- Beside it, an elegant lady's \VUgras{writing desk} \textsuperscript{(5)}
stands open. It is where the \VUgras{lady of the house} \textsuperscript{(23)}
writes her letters.

%%K t06-c2-04-1 p
4. --- The window is dressed with \VUgras{net curtains} \textsuperscript{(6)} held
back by \VUgras{tiebacks} \textsuperscript{(8)} hooked to the \VUgras{pegs}
\textsuperscript{(7)}, and with heavy fabric curtains topped by a
\VUgras{valance} \textsuperscript{(9)}.

%%K t06-c2-05-1 p
5. --- In the window recess, a \VUgras{jardinière} \textsuperscript{(11)} holds a
green \VUgras{plant} \textsuperscript{(10)}, a magnificent palm whose leaves reach
almost as far as the \VUgras{oil lamp} \textsuperscript{(12)}.

%%K t06-c2-06-1 p
6. --- The \VUgras{lampshade} \textsuperscript{(13)} was made by Mary, the charming
\VUgras{young lady} \textsuperscript{(14)} at the \VUgras{piano} \textsuperscript{(15)}.
Sitting very straight on the \VUgras{stool} \textsuperscript{(16)}, she is
practising a piece. She is very skilful and very tidy, as her
\VUgras{music cabinet} \textsuperscript{(17)} shows: it is in perfect order.

%%K t06-tit-4 sub
\VUsaut{5.0mm}
\VUpk{10.2pt}{40}{The Rascal.}
\VUsaut{3.0mm}

%%K t06-c2-07-1 p
7. --- Her brother Paul is looking for a \VUgras{book} \textsuperscript{(19)} in the
\VUgras{bookcase} \textsuperscript{(18)}. He is a \VUgras{boy} \textsuperscript{(20)} who is
restless and quite unbearable. Hanging on to the bookcase to reach the
book that is too high up, he is in danger of bringing down the
\VUgras{bust} \textsuperscript{(21)} on top of it.

%%K t06-c2-08-1 p
8. --- He is taking advantage, to do this silly thing, of his mother
being busy talking with a friend, a \VUgras{lady} \textsuperscript{(22)} who has come
to spend a few days with her. The mother is sitting on the \VUgras{sofa}
\textsuperscript{(24)} beside the \VUgras{armchair} \textsuperscript{(25)}, and her friend
is on the \VUgras{chair} \textsuperscript{(27)}, of which we can see only the
\VUgras{back} \textsuperscript{(26)}.

%%K t06-c2-09-1 p
9. --- The big \VUgras{frame} \textsuperscript{(30)} hanging on the wall holds the
\VUgras{portrait} \textsuperscript{(29)} of a relative. In the \VUgras{album}
\textsuperscript{(32)} lying on the table are gathered the \VUgras{photographs}
\textsuperscript{(33)} of the rest of the family. The children were leafing
through it a moment ago, for the \VUgras{chairs} \textsuperscript{(31)} are still
grouped around the table.

%%K t06-c2-10-1 p
10. --- The porcelain \VUgras{Chinese vase} \textsuperscript{(34)} near the
\VUgras{paper knife} \textsuperscript{(35)} was given to Mary on her name day, and
the \VUgras{folding screen} \textsuperscript{(36)} that fills the corner of the room
was brought back from China by her uncle.

%%K t06-c2-11-1 p
11. --- Brightened by the handsome \VUgras{pictures} \textsuperscript{(37)} hanging
on the \VUgras{wall} \textsuperscript{(42)}, lit by the ceiling
\VUgras{chandelier} \textsuperscript{(38)}, whose \VUgras{electric lamps}
\textsuperscript{(39)} shine so cheerfully in the evening, this drawing room
looks very pleasant. The soft \VUgras{carpet} \textsuperscript{(40)} spread over
the floor, and the \VUgras{electric bell} \textsuperscript{(41)} that is so handy,
make it perfectly comfortable.

%%K t06-tit-5 sub
\VUsaut{3.6mm}
\VUpk{10.2pt}{40}{The Dining Room.}
\VUsaut{3.0mm}

%%K t06-c3-01-1 p
1. --- The dinner bell has rung. The whole family, except Mary, is
gathered around the table. The dining room is pleasantly warmed by an
excellent earthenware \VUgras{stove} \textsuperscript{(1)} with a slender
\VUgras{pipe} \textsuperscript{(2)}.

%%K t06-c3-02-1 p
2. --- This room does not hold much furniture. Along the wall, above the
\VUgras{stool} \textsuperscript{(4)}, hangs a decorative little \VUgras{fountain}
\textsuperscript{(3)}. Close by is the \VUgras{sideboard} \textsuperscript{(5)}, on which
the cold dishes, the desserts and the various pieces of tableware not
needed at the moment are set down.

%%K t06-c3-03-1 p
3. --- So on the upper shelf we can see a \VUgras{roast chicken}
\textsuperscript{(7)}, a \VUgras{fish} \textsuperscript{(8)} and a \VUgras{bowl}
\textsuperscript{(10)} of \VUgras{fruit} \textsuperscript{(9)}. Below are the
\VUgras{cheese} \textsuperscript{(11)}, the \VUgras{bread} \textsuperscript{(12)} and the
\VUgras{cruet stand} \textsuperscript{(13)}, which holds everything needed to dress
the salad: the oil, the vinegar, the \VUgras{salt} \textsuperscript{(14)} and the
\VUgras{pepper} \textsuperscript{(15)}. Lastly, on the bottom shelf, there is a
\VUgras{cake} \textsuperscript{(17)} beside the \VUgras{mustard pot} \textsuperscript{(16)}.

%%K t06-c3-04-1 p
4. --- The dining room clock, which is called a \VUgras{wall clock}
\textsuperscript{(20)}, is a very unusual shape. It is set in a wooden frame that
also holds a \VUgras{barometer} \textsuperscript{(19)} and a \VUgras{thermometer}
\textsuperscript{(18)}.

%%K t06-tit-6 sub
\VUsaut{4.6mm}
\VUpk{10.2pt}{40}{Serving Dinner.}
\VUsaut{3.0mm}

%%K t06-c3-05-1 p
5. --- All round the room the walls are lined with \VUgras{wooden panelling}
\textsuperscript{(21)} of waxed oak. A fabric \VUgras{door curtain} \textsuperscript{(22)}
closes off well enough the door leading to the veranda. That is where the
family will go to take coffee after dinner. The \VUgras{cups}
\textsuperscript{(24)} and \VUgras{saucers} \textsuperscript{(25)} are already set out on
the \VUgras{china cabinet} \textsuperscript{(23)}.

%%K t06-c3-06-1 p
6. --- Above the table a \VUgras{hanging lamp} \textsuperscript{(26)} is fixed to
the ceiling; its very bright light fills the whole dining room in the
evening.

%%K t06-c3-07-1 p
7. --- Let us now look at the \VUgras{dining table} \textsuperscript{(27)} itself.
When the family came in, the table was laid. On the \VUgras{tablecloth}
\textsuperscript{(28)} there had been set, at each diner's place, a \VUgras{plate}
\textsuperscript{(33)}, a \VUgras{napkin} \textsuperscript{(29)}, a \VUgras{spoon}
\textsuperscript{(34)}, a \VUgras{fork} \textsuperscript{(35)}, a \VUgras{knife}
\textsuperscript{(36)} and a \VUgras{glass} \textsuperscript{(32)}. There were also
\VUgras{bottles} \textsuperscript{(30)} for the wine and a \VUgras{carafe}
\textsuperscript{(31)} for the water.

%%K t06-c3-08-1 p
8. --- The \VUgras{servant} \textsuperscript{(39)} first brought in the
\VUgras{soup tureen} \textsuperscript{(37)}, in which some soup is still steaming.
He is now serving the first \VUgras{course} \textsuperscript{(38)}, a meat dish.
Then he will hand round the ones we have already named; and finally,
after the \VUgras{dessert}, he will take advantage of his employers having
gone through to the veranda to clear the table and put the dining room
straight again.

%%K t06-c3-08-2 p
\textit{Note.} --- \textit{The everyday words to do with food are given
here only in outline. The list will be completed in the two charts “The
Hotel” (no. 13) and “The Market” (no. 15).}

%%K t06-tit-7 sub
\VUsaut{4.6mm}
\VUpk{10.2pt}{40}{The Kitchen.}
\VUsaut{3.0mm}

%%K t06-c4-01-1 p
1. --- In the kitchen, at which we glance through the half-open door, we
find a picturesque muddle. In front of the \VUgras{hearth} \textsuperscript{(1)},
where the \VUgras{fire} \textsuperscript{(3)} crackles, the \VUgras{roasting jack}
\textsuperscript{(5)} is set up. A fine \VUgras{roast} \textsuperscript{(7)} is on the
\VUgras{spit} \textsuperscript{(6)} and is finishing cooking.

%%K t06-c4-02-1 p
2. --- The \VUgras{bellows} \textsuperscript{(8)} and the \VUgras{coal shovel}
\textsuperscript{(9)}, which have just been used, have been left on the floor,
along with the \VUgras{logs} \textsuperscript{(10)}.

%%K t06-c4-03-1 p
3. --- The top of the mantelpiece is tidy: the \VUgras{lantern}
\textsuperscript{(11)}, the \VUgras{matchbox holder} \textsuperscript{(13)} with the
\VUgras{matches} \textsuperscript{(12)} in it, the \VUgras{candlestick}
\textsuperscript{(14)} and the \VUgras{chamber candlestick} \textsuperscript{(15)} with
its \VUgras{candle} \textsuperscript{(16)} are all in their places. But the big
\VUgras{coal scuttle} \textsuperscript{(18)}, full of the \VUgras{coal}
\textsuperscript{(17)} the \VUgras{cook} \textsuperscript{(23)} has just been down to the
cellar to fetch, has been left in the middle of the kitchen.

%%K t06-c4-04-1 p
4. --- The truth is that the poor woman has to hurry if lunch is to be
ready on time. Just now she is at the \VUgras{cooking range} \textsuperscript{(19)},
stirring a \VUgras{sauce} \textsuperscript{(24)} that must not be allowed to burn.

%%K t06-c4-05-1 p
5. --- In the \VUgras{oven} \textsuperscript{(20)} a cake is baking, which she has
made for the children's tea. Above the range hang the \VUgras{ladle}
\textsuperscript{(21)} and the \VUgras{skimmer} \textsuperscript{(22)}.

%%K t06-tit-8 sub
\VUsaut{4.0mm}
\VUpk{10.2pt}{40}{The Kitchenware.}
\VUsaut{3.0mm}

%%K t06-c4-06-1 p
6. --- The \VUgras{set of pans} \textsuperscript{(25)} hanging at the back of the
room is very clean. The handsome copper \VUgras{saucepans} \textsuperscript{(26)},
the \VUgras{milk jug} \textsuperscript{(27)}, the \VUgras{strainer} \textsuperscript{(28)}
and the \VUgras{grater} \textsuperscript{(29)} have been carefully scoured.

%%K t06-c4-07-1 p
7. --- The \VUgras{salt box} \textsuperscript{(30)} stands on the small dresser
beside the \VUgras{teapot} \textsuperscript{(31)} and the \VUgras{oil demijohn}
\textsuperscript{(32)}. The \VUgras{drawer} \textsuperscript{(33)} probably holds the
cutlery and the kitchen knives.

%%K t06-c4-08-1 p
8. --- The \VUgras{broom} \textsuperscript{(34)} and the \VUgras{feather duster}
\textsuperscript{(35)} are in a corner. The cook has just been using the
\VUgras{chopping block} \textsuperscript{(36)}; she has left it near the window.

%%K t06-c4-09-1 p
9. --- A \VUgras{hand towel} \textsuperscript{(37)} hangs on the wall next to the
\VUgras{funnel} \textsuperscript{(38)}. In this part of the kitchen are kept the
\VUgras{grill} \textsuperscript{(39)}, the \VUgras{chopper} \textsuperscript{(40)} and the
\VUgras{chopping board} \textsuperscript{(41)}. Then, above the chopper, on a
\VUgras{shelf} \textsuperscript{(42)}, are \VUgras{pots} \textsuperscript{(43)}, the
\VUgras{stew pot} \textsuperscript{(44)} with its \VUgras{lid} \textsuperscript{(45)}, and
the \VUgras{cauldron} \textsuperscript{(46)}. The \VUgras{frying pan} \textsuperscript{(47)}
hangs close by.

%%K t06-c4-10-1 p
10. --- Only the copper \VUgras{tap} \textsuperscript{(48)} throws a gleam into this
corner. The \VUgras{sink} \textsuperscript{(49)}, on which the big \VUgras{pitcher}
\textsuperscript{(50)} stands, must be very convenient with its little cupboard,
where the \VUgras{vinegar cask} \textsuperscript{(51)} with its \VUgras{tap}
\textsuperscript{(52)}, the \VUgras{rubbish box} \textsuperscript{(53)} and the
\VUgras{dirty dishes} \textsuperscript{(54)} are kept.

%%K t06-tit-9 sub
\VUsaut{4.0mm}
\VUpk{10.2pt}{40}{Other things kept in the kitchen.}
\VUsaut{3.0mm}

%%K t06-c4-11-1 p
11. --- The big \VUgras{tub} \textsuperscript{(55)} in which the \VUgras{vegetables}
\textsuperscript{(58)} are washed, and the cast-iron \VUgras{cauldron}
\textsuperscript{(56)} standing on the \VUgras{tiled floor} \textsuperscript{(57)},
probably belong in that cupboard too.

%%K t06-c4-12-1 p
12. --- The cook is certainly not short of stoves, for here on the corner
of the table is a \VUgras{gas ring} \textsuperscript{(59)} with water heating on it
in a small saucepan. The \VUgras{steam} \textsuperscript{(60)} escaping from it
tells us it must be boiling. This water is probably for the coffee, for
here at the other end of the table are the \VUgras{coffee pot}
\textsuperscript{(64)} and the \VUgras{coffee grinder} \textsuperscript{(63)}.

%%K t06-c4-13-1 p
13. --- The gas ring is connected to the \VUgras{gas burner} \textsuperscript{(62)}
by a long \VUgras{rubber tube} \textsuperscript{(61)}.

%%K t06-c4-14-1 p
14. --- The \VUgras{daily help} \textsuperscript{(66)}, who comes to assist the
cook, is preparing the vegetables for dinner. When they have been peeled
and washed, she will put them in the \VUgras{pan} \textsuperscript{(65)} until it
is time to cook them.

%%K t06-tit-10 sub
\VUsaut{4.0mm}
{\centering\fontsize{10.2pt}{10.2pt}\selectfont\textls[40]{\scshape The Bathroom.} \textit{(The bathing room.)}\par}
\VUsaut{3.0mm}

%%K t06-c5-01-1 p
1. --- Only one more piece of prying remains before we know the main
rooms of the house: to look at the fittings of the bathroom. It will not
take long, for it is not large, and we shall quickly see that the
\VUgras{bathtub} \textsuperscript{(1)} has a good \VUgras{water heater}
\textsuperscript{(2)}, that the \VUgras{soap dish} \textsuperscript{(3)} is within reach
of the bather's hand, and that the \VUgras{towel rail} \textsuperscript{(5)} must
make it easy to dry the \VUgras{towels} \textsuperscript{(4)}.

%%K t06-c5-02-1 p
2. --- This room has its walls covered with \VUgras{glazed tiles}
\textsuperscript{(6)}, which made it possible to install a \VUgras{shower}
\textsuperscript{(7)} without any fear that the water splashing about while it is
used might damage them. A small zinc \VUgras{basin} \textsuperscript{(8)} for
footbaths completes the furnishings.
\VUsaut{6.0mm}
\VUfilet{22mm}
